\documentclass[10pt]{article}
\usepackage[margin=1in, paperwidth=8.5in, paperheight=11in]{geometry}
\usepackage{ifpdf,amsmath, amssymb, comment, color, graphicx, stmaryrd,setspace,enumitem,tikz, fancyhdr, wrapfig, textcomp, units, mathptmx, pdfpages, cancel, multicol}
\usepackage{hyperref}
\hypersetup{
    colorlinks=true,
    urlcolor=blue,
}

\setlength{\headheight}{14.5pt}
\newcommand{\Q}{\mathbb{Q}}
\newcommand{\R}{\mathbb{R}}
\newcommand{\Z}{\mathbb{Z}}
\newcommand{\vu}{\mathbf{u}}
\newcommand{\vv}{\mathbf{v}}
\newcommand{\vw}{\mathbf{w}}
\newcommand{\vi}{\mathbf{i}}
\newcommand{\vj}{\mathbf{j}}
\newcommand{\vk}{\mathbf{k}}
\newcommand{\vn}{\mathbf{n}}
\newcommand{\vr}{\mathbf{r}}
\newcommand{\va}{\mathbf{a}}
\newcommand{\vF}{\mathbf{F}}
\newcommand{\vL}{\mathbf{L}}
\newcommand{\vT}{\mathbf{T}}
\newcommand{\vN}{\mathbf{N}}
\newcommand{\proj}{\operatorname{proj}}
\newcommand{\orth}{\operatorname{orth}}
\newcommand\del\nabla
\newcommand\dotp[1][.5]{\,\mathbin{\vcenter{\hbox{\scalebox{#1}{$\bullet$}}}}\,}

% Solution text is in red. If you want the solutions to show, remove the \iffalse from the definition of the \red command.
%\newcommand{\red}[1]{ %\iffalse
%	\textcolor{red}{#1} }%\fi}

\newenvironment{red}{\color{red}}{\ignorespacesafterend}
\newenvironment{blue}{\color{blue}}{\ignorespacesafterend}
\newcommand{\green}[1]{\textcolor{green}{#1}}
\renewcommand{\section}[1]{\begin{center} \textbf{#1} \\\end{center}}
%
\hyphenpenalty=5000
\setlength{\parindent}{0in}
%\oddsidemargin=-.25in
\allowdisplaybreaks
\pagestyle{fancy}
\renewcommand{\headrulewidth}{0pt}
\lhead{MATH 202}
\rhead{Spring 2022}
%\lfoot{\copyright\ CLEAR Calculus 2010}
\cfoot{}

\begin{document}
%


%\onehalfspacing
\allowdisplaybreaks
%##################################################################
\section{PS\#4: Checkpoint Preview M1-M5 - \red{Answer key} }

Because many of these problems are randomly generated, I can't write meaningful answer keys beyond the ``partial answers'' provided by Canvas. However, everyone got the same \#1, so here's an answer key for that.

\everymath{\displaystyle}
\begin{enumerate}[leftmargin=0pt]

\item Consider the function $f(x)=e^{-x^2}$, which does not have an elementary antiderivative. Define the integral function $F(t) = \int_0^t f(x)\, dx$.

\begin{enumerate}
    \item What does the Second FTC tell us about the relationship between $F$ and $f$?
    
    \begin{red}
    The Second FTC says that $F$ is an antiderivative of $f$ -- that is, that $F' = f$.
    \end{red}
    \item Use appropriate technology to approximate $F(2)$ using left, right, midpoint, and trapezoid sums with 10 subintervals. Explain which of these approximations are overestimates and which are underestimates.
    
    \begin{red}
    I used the following commands in Geogebra to produce the following estimates. You can certainly use different technology as long as it produces the same calculations. After defining $f(x) = e^{-x^2}$:
    
    \begin{itemize}
        \item Left sum: \texttt{RectangleSum(f(x),0,2,10,0)} $\approx 0.9800072469099$
        \item Right sum: \texttt{RectangleSum(f(x),0,2,10,1)} $\approx 0.7836703746876$ 
        \item Midpoint sum: \texttt{RectangleSum(f(x),0,2,10,0.5)} $\approx 0.8822020699923$
        \item Trapezoid sum: \texttt{TrapezoidalSum(f(x),0,2,10)} $\approx 0.8818388107988$
    \end{itemize}
    \end{red}
    \item Use the results from part B to approximate $F(2)$ using Simpson's rule.
    
    \begin{red}
    Simpson's rule is the weighted average of the midpoint sum and the trapezoid sum:
    
    \begin{align*}
        S_{20} &= \frac{2M_n + T_n}{3} 
        = \frac{1}{3}(2 \cdot 0.8822020699923 + 0.8818388107988) \\
        &= \frac{1}{3}( 2.64624295078 ) = 0.88208098359
    \end{align*}
    \end{red}
\end{enumerate}

\end{enumerate}
\end{document}